\section{Kết quả và thảo luận}
\label{sec:results}

\subsection{So sánh tổng thể trên cấu hình 1 CPU--512 MB}

Bảng~\ref{tab:overall_results} trình bày kết quả chính trên profile \texttt{edge\_1cpu\_512mb}. Với \method{M5} và \method{M6}, bảng chọn FAR budget $\alpha=0{,}02$ làm điểm vận hành đại diện vì đây là mức rủi ro thấp nhưng không quá cực đoan. Các số liệu được lấy từ \texttt{combined\_edge\_summary.csv}.

\begin{table}[!htbp]
\centering
\caption{So sánh tổng thể trên profile \texttt{edge\_1cpu\_512mb}. Với \method{M5}/\method{M6}, $\alpha=0{,}02$.}
\label{tab:overall_results}
\resizebox{\textwidth}{!}{
\begin{tabular}{lcccccc}
\toprule
\textbf{Method} & \textbf{FAR\_active} & \textbf{FRR\_active} & \textbf{FRR + defer} & \textbf{Defer} & \textbf{Automation} & \textbf{P95 latency / RAM} \\
\midrule
\method{M0} Always-fast fixed        & 5,0\% & 0,0\% & 3,2\% & 3,2\% & 96,8\% & 483,9 ms / 219,4 MB \\
\method{M1} Always-fast bin          & 1,7\% & 0,0\% & 3,2\% & 3,2\% & 96,8\% & 484,4 ms / 219,1 MB \\
\method{M4} Conditional path bin     & 1,6\% & 0,0\% & 0,0\% & 0,0\% & 100,0\% & 484,0 ms / 219,2 MB \\
\method{M5} Conditional risk         & 0,0\% & 1,6\% & 1,6\% & 0,0\% & 100,0\% & 491,7 ms / 219,5 MB \\
\method{M6} Conditional risk + defer & 0,0\% & 0,0\% & 1,6\% & 0,8\% & 99,2\% & 482,8 ms / 211,9 MB \\
\bottomrule
\end{tabular}
}
\end{table}

Kết quả cho thấy \method{M0} có FAR khoảng 5,0\%, vượt các FAR budget thấp như 0,01--0,03. \method{M4} giảm FAR xuống khoảng 1,6\% nhờ chọn path theo điều kiện và dùng ngưỡng theo bin, nhưng phương pháp này chưa trực tiếp hiệu chỉnh theo FAR budget. \method{M5} đặt ràng buộc risk-constrained rõ ràng: tại $\alpha=0{,}02$, test FAR\_active bằng 0,0\%, đổi lại FRR\_active là 1,6\%.

\method{M6} là phương pháp tốt nhất trong benchmark hiện tại: test FAR\_active 0,0\%, FRR\_active 0,0\%, defer rate 0,8\% và automation rate 99,2\%. Điều này không có nghĩa hệ thống ``nhận đúng tất cả''; nếu tính genuine bị defer như một trường hợp chưa phục vụ thành công, FRR + defer là 1,6\%. Cách báo cáo này minh bạch hơn vì defer được tách riêng thay vì bị ẩn trong accuracy.

\subsection{Kết quả theo condition bin tại FAR budget 0,02}

Bảng~\ref{tab:condition_results} so sánh \method{M5} và \method{M6} theo từng condition bin ở FAR budget 0,02. Các số liệu được lấy từ \texttt{combined\_far\_budget\_summary.csv}.

\begin{table}[!htbp]
\centering
\caption{So sánh theo condition bin trên \texttt{edge\_1cpu\_512mb}, FAR budget $\alpha=0{,}02$.}
\label{tab:condition_results}
\begin{tabular}{llccccc}
\toprule
\textbf{Method} & \textbf{Condition} & \textbf{FAR\_active} & \textbf{FRR\_active} & \textbf{FRR + defer} & \textbf{Defer} & \textbf{Automation} \\
\midrule
\method{M5} & bright  & 0,0\% & 5,0\% & 5,0\% & 0,0\% & 100,0\% \\
\method{M5} & medium  & 0,0\% & 0,0\% & 0,0\% & 0,0\% & 100,0\% \\
\method{M5} & dark    & 0,0\% & 0,0\% & 0,0\% & 0,0\% & 100,0\% \\
\method{M5} & overall & 0,0\% & 1,6\% & 1,6\% & 0,0\% & 100,0\% \\
\midrule
\method{M6} & bright  & 0,0\% & 0,0\% & 5,0\% & 2,5\% & 97,5\% \\
\method{M6} & medium  & 0,0\% & 0,0\% & 0,0\% & 0,0\% & 100,0\% \\
\method{M6} & dark    & 0,0\% & 0,0\% & 0,0\% & 0,0\% & 100,0\% \\
\method{M6} & overall & 0,0\% & 0,0\% & 1,6\% & 0,8\% & 99,2\% \\
\bottomrule
\end{tabular}
\end{table}

Trong test split hiện tại, bin \texttt{dark} đạt FAR\_active 0,0\% và FRR\_active 0,0\% cho cả \method{M5} và \method{M6} tại $\alpha=0{,}02$. Vì số cặp test nhỏ, không nên kết luận rằng điều kiện tối đã được giải quyết hoàn toàn. Kết luận hợp lý hơn là: dưới tập test hiện tại, phương pháp risk-constrained không làm tăng FAR trong bin tối, và defer giúp loại bỏ lỗi active còn lại ở mức overall với chi phí tự động hóa rất nhỏ.

\subsection{Ngưỡng calibration của M6}

Bảng~\ref{tab:thresholds} trình bày ngưỡng accept/reject của \method{M6} tại FAR budget 0,02 trên profile \texttt{edge\_1cpu\_512mb}. Các số liệu được lấy từ \texttt{combined\_calibration\_thresholds.csv}.

\begin{table}[!htbp]
\centering
\caption{Ngưỡng calibration của \method{M6} tại $\alpha=0{,}02$.}
\label{tab:thresholds}
\begin{tabular}{lcccc}
\toprule
\textbf{Condition} & \textbf{$\tau_{accept}$} & \textbf{$\tau_{reject}$} & \textbf{Calib. FAR} & \textbf{Calib. FRR} \\
\midrule
\texttt{bright} & 0,601 & 0,571 & 0,0\% & 0,0\% \\
\texttt{medium} & 0,871 & 0,841 & 0,0\% & 0,0\% \\
\texttt{dark} & 0,554 & 0,524 & 0,0\% & 0,0\% \\
\bottomrule
\end{tabular}
\end{table}

Bảng này cho thấy ngưỡng của bin \texttt{dark} thấp hơn \texttt{bright}, nhưng điểm quan trọng là ngưỡng được chọn qua ràng buộc FAR trên calibration, không phải hạ thủ công vì trời tối. Với \method{M6}, score nằm giữa $\tau_{reject}$ và $\tau_{accept}$ sẽ bị defer, do đó vùng rủi ro không bị tự động chấp nhận.

\subsection{So sánh các cấu hình edge-constrained}

Bảng~\ref{tab:profile_results} báo cáo latency và RAM của \method{M6} tại FAR budget 0,02 trên ba cấu hình AWS. Các số liệu được lấy từ \texttt{combined\_edge\_summary.csv}.

\begin{table}[!htbp]
\centering
\caption{Chi phí tài nguyên của \method{M6} tại $\alpha=0{,}02$ trên ba profile AWS.}
\label{tab:profile_results}
\begin{tabular}{lccccc}
\toprule
\textbf{Profile} & \textbf{CPU} & \textbf{Memory} & \textbf{Mean latency} & \textbf{P95 latency} & \textbf{RAM peak} \\
\midrule
\texttt{edge\_1cpu\_512mb} & 1 & 512 MB & 430,4 ms & 482,8 ms & 211,9 MB \\
\texttt{edge\_1cpu\_1gb}   & 1 & 1 GB   & 438,4 ms & 492,9 ms & 221,5 MB \\
\texttt{edge\_2cpu\_2gb}   & 2 & 2 GB   & 324,2 ms & 362,5 ms & 210,9 MB \\
\bottomrule
\end{tabular}
\end{table}

Khi tăng từ 1 CPU lên 2 CPU, P95 latency giảm từ khoảng 483--493 ms xuống khoảng 363 ms. RAM peak nằm quanh 211--222 MB trong Docker, thấp hơn giới hạn 512 MB của profile chặt nhất. Tuy nhiên, đây vẫn là kết quả trên AWS ARM có Docker limit, không phải số đo trực tiếp trên board biên vật lý.

\subsection{Diễn giải và giới hạn của kết quả}

Kết quả ủng hộ hướng risk-constrained theo ba điểm. Thứ nhất, so sánh tại cùng FAR budget rõ ràng hơn so với chỉ báo cáo FRR/FAR rời rạc. Thứ hai, defer nhỏ có thể giảm lỗi active mà vẫn giữ automation rate cao. Thứ ba, conditional path dùng robust path khoảng 21\% số mẫu, nghĩa là hệ thống không luôn bật pipeline nặng.

Dù vậy, có ba giới hạn cần nhấn mạnh. Tập test chỉ có 124 cặp nên một vài phần trăm tương ứng với rất ít mẫu; vì vậy FAR 0,0\% không nên được hiểu là bảo đảm an toàn tuyệt đối. Ngoài ra, tại FAR budget 0,05, \method{M5} có test FAR 6,5\%, cao hơn mức calibration mong muốn, cho thấy sai lệch giữa calibration/test có thể xảy ra khi dữ liệu nhỏ. Cuối cùng, dữ liệu tối và nhiễu cần được kiểm chứng thêm bằng ảnh thật thu từ camera biên trong điều kiện ngày/đêm tự nhiên.
